% Importance sampling for a rare-event tail. The quantity of interest is
% Pr(X > t) for a standard normal X at a far threshold t. Naive Monte
% Carlo draws from N(0,1) and almost never lands past t, so the estimator
% has enormous relative variance. Importance sampling draws from a shifted
% proposal centred near t (here N(t,1)) and reweights by the likelihood
% ratio, placing samples where the integrand lives. The figure overlays
% the target density, the tail region, and the shifted proposal.
\begin{figure}[htbp]
\centering
\begin{tikzpicture}
\begin{axis}[
  width=11cm, height=6.4cm,
  xlabel={$x$},
  ylabel={density},
  xmin=-3.5, xmax=6.5, ymin=0, ymax=0.45,
  domain=-3.5:6.5, samples=180,
  legend pos=north east,
  legend cell align=left,
  grid=both, grid style={enggray!18},
  tick label style={font=\small},
  label style={font=\small},
  legend style={font=\scriptsize},
]

% Target density N(0,1)
\addplot[very thick, engblue, domain=-3.5:6.5, samples=180]
  {0.3989422804*exp(-0.5*x*x)};
\addlegendentry{target $f=\mathcal{N}(0,1)$}

% Shifted proposal N(3,1)
\addplot[very thick, engred, domain=-3.5:6.5, samples=180]
  {0.3989422804*exp(-0.5*(x-3)*(x-3))};
\addlegendentry{proposal $g=\mathcal{N}(3,1)$}

% Shade the tail region x > 3 under the target density
\addplot[fill=engblue!22, draw=none, domain=3:6.5, samples=60]
  {0.3989422804*exp(-0.5*x*x)} \closedcycle;

% threshold marker
\draw[dashed, engblack] (axis cs:3,0) -- (axis cs:3,0.4);
\node[anchor=south, engblack, font=\scriptsize] at (axis cs:3,0.4)
  {threshold $t=3$};

\end{axis}
\end{tikzpicture}
\caption[Importance sampling for a rare tail]{Estimating a rare-tail
probability $\Pr(X>t)$ by naive Monte Carlo wastes almost every draw:
samples from the target $\mathcal{N}(0,1)$ (blue) land in the shaded tail
past $t=3$ only about once in $740$ draws, so the estimator's relative
error is large. Importance sampling instead draws from a proposal (red)
shifted to sit over the tail, then corrects each draw by the likelihood
ratio $f(x)/g(x)$. Every sample now contributes to the estimate, and the
variance drops by orders of magnitude for the same number of draws. The
method is the standard tool when the event of interest is rare and a
plain average would need an impractical sample size. Source: schematic by
the authors; the estimator and its variance are worked in section 13.6.}
\label{fig:vol02ch13:importance-sampling}
\end{figure}
